% sos_decomposition.tex — Parrilo's SOS decomposition, in this repository's
% notation, with one fully worked example.
%
% Audience: an advisor skimming the GitHub profile who wants to know in ninety
% seconds whether the author understands SOS or has only called a library. Write
% for that reader. The worked example in Section 4 is what answers the question;
% everything else is setup.
%
% Compile:  latexmk -pdf sos_decomposition.tex   (from docs/derivations/)

\documentclass[11pt]{article}
% preamble.tex — shared macros for the derivations.
%
% Both derivation files are standalone-compilable via % preamble.tex — shared macros for the derivations.
%
% Both derivation files are standalone-compilable via % preamble.tex — shared macros for the derivations.
%
% Both derivation files are standalone-compilable via \input{preamble}. Compile
% from inside docs/derivations/ so the relative \input resolves:
%
%   latexmk -pdf sos_decomposition.tex
%
% Notation here must match docs/theory_notes.md §1 and the code identifiers.
% If they drift, the cross-references in the notebooks stop meaning anything.

\usepackage[utf8]{inputenc}
\usepackage[T1]{fontenc}
\usepackage{amsmath, amssymb, amsthm}
\usepackage{mathtools}
\usepackage{bm}
\usepackage{graphicx}
\usepackage{booktabs}
\usepackage[margin=1in]{geometry}
\usepackage[colorlinks=true, linkcolor=blue, citecolor=blue]{hyperref}
\usepackage{cleveref}   % requires two pdflatex passes

\theoremstyle{plain}
\newtheorem{theorem}{Theorem}
\newtheorem{lemma}[theorem]{Lemma}
\newtheorem{proposition}[theorem]{Proposition}
\newtheorem{corollary}[theorem]{Corollary}

\theoremstyle{definition}
\newtheorem{definition}[theorem]{Definition}
\newtheorem{example}[theorem]{Example}
\newtheorem{assumption}[theorem]{Assumption}

\theoremstyle{remark}
\newtheorem{remark}[theorem]{Remark}

% --- sets and cones ----------------------------------------------------------
\newcommand{\R}{\mathbb{R}}
\newcommand{\N}{\mathbb{N}}
\newcommand{\Sym}{\mathbb{S}}                      % symmetric matrices
\newcommand{\PSD}{\succeq 0}
\newcommand{\SOS}{\Sigma}                          % the SOS cone, \SOS[x]
\newcommand{\Rx}{\R[x]}                            % polynomial ring
\newcommand{\Mset}{\mathcal{M}}                    % quadratic module

% --- the objects of the certificate ------------------------------------------
\newcommand{\Vdot}{\dot{V}}
\newcommand{\Lie}[2]{\nabla #1 \cdot #2}           % \Lie{V}{f}
\newcommand{\sublevel}[2]{\Omega_{#2}}             % \sublevel{V}{\rho}
\newcommand{\ROA}{\mathcal{R}}                     % region of attraction
\newcommand{\gram}{Q}
\newcommand{\basis}{z}

% --- operators ---------------------------------------------------------------
\DeclareMathOperator{\tr}{tr}
\DeclareMathOperator{\rank}{rank}
\DeclareMathOperator{\diag}{diag}
\DeclareMathOperator{\deg}{deg}
\DeclareMathOperator*{\argmax}{arg\,max}
\DeclareMathOperator*{\argmin}{arg\,min}

% --- a macro for pointing prose at the implementation -------------------------
% Use \codelink{lie_derivative}{julia/src/polynomial_utils.jl} in the text every
% time a symbol first appears. A derivation that names the function implementing
% it is checkable; one that does not is a separate document that happens to sit
% in the same repository.
\newcommand{\codelink}[2]{\texttt{#1} (\texttt{#2})}
. Compile
% from inside docs/derivations/ so the relative \input resolves:
%
%   latexmk -pdf sos_decomposition.tex
%
% Notation here must match docs/theory_notes.md §1 and the code identifiers.
% If they drift, the cross-references in the notebooks stop meaning anything.

\usepackage[utf8]{inputenc}
\usepackage[T1]{fontenc}
\usepackage{amsmath, amssymb, amsthm}
\usepackage{mathtools}
\usepackage{bm}
\usepackage{graphicx}
\usepackage{booktabs}
\usepackage[margin=1in]{geometry}
\usepackage[colorlinks=true, linkcolor=blue, citecolor=blue]{hyperref}
\usepackage{cleveref}   % requires two pdflatex passes

\theoremstyle{plain}
\newtheorem{theorem}{Theorem}
\newtheorem{lemma}[theorem]{Lemma}
\newtheorem{proposition}[theorem]{Proposition}
\newtheorem{corollary}[theorem]{Corollary}

\theoremstyle{definition}
\newtheorem{definition}[theorem]{Definition}
\newtheorem{example}[theorem]{Example}
\newtheorem{assumption}[theorem]{Assumption}

\theoremstyle{remark}
\newtheorem{remark}[theorem]{Remark}

% --- sets and cones ----------------------------------------------------------
\newcommand{\R}{\mathbb{R}}
\newcommand{\N}{\mathbb{N}}
\newcommand{\Sym}{\mathbb{S}}                      % symmetric matrices
\newcommand{\PSD}{\succeq 0}
\newcommand{\SOS}{\Sigma}                          % the SOS cone, \SOS[x]
\newcommand{\Rx}{\R[x]}                            % polynomial ring
\newcommand{\Mset}{\mathcal{M}}                    % quadratic module

% --- the objects of the certificate ------------------------------------------
\newcommand{\Vdot}{\dot{V}}
\newcommand{\Lie}[2]{\nabla #1 \cdot #2}           % \Lie{V}{f}
\newcommand{\sublevel}[2]{\Omega_{#2}}             % \sublevel{V}{\rho}
\newcommand{\ROA}{\mathcal{R}}                     % region of attraction
\newcommand{\gram}{Q}
\newcommand{\basis}{z}

% --- operators ---------------------------------------------------------------
\DeclareMathOperator{\tr}{tr}
\DeclareMathOperator{\rank}{rank}
\DeclareMathOperator{\diag}{diag}
\DeclareMathOperator{\deg}{deg}
\DeclareMathOperator*{\argmax}{arg\,max}
\DeclareMathOperator*{\argmin}{arg\,min}

% --- a macro for pointing prose at the implementation -------------------------
% Use \codelink{lie_derivative}{julia/src/polynomial_utils.jl} in the text every
% time a symbol first appears. A derivation that names the function implementing
% it is checkable; one that does not is a separate document that happens to sit
% in the same repository.
\newcommand{\codelink}[2]{\texttt{#1} (\texttt{#2})}
. Compile
% from inside docs/derivations/ so the relative \input resolves:
%
%   latexmk -pdf sos_decomposition.tex
%
% Notation here must match docs/theory_notes.md §1 and the code identifiers.
% If they drift, the cross-references in the notebooks stop meaning anything.

\usepackage[utf8]{inputenc}
\usepackage[T1]{fontenc}
\usepackage{amsmath, amssymb, amsthm}
\usepackage{mathtools}
\usepackage{bm}
\usepackage{graphicx}
\usepackage{booktabs}
\usepackage[margin=1in]{geometry}
\usepackage[colorlinks=true, linkcolor=blue, citecolor=blue]{hyperref}
\usepackage{cleveref}   % requires two pdflatex passes

\theoremstyle{plain}
\newtheorem{theorem}{Theorem}
\newtheorem{lemma}[theorem]{Lemma}
\newtheorem{proposition}[theorem]{Proposition}
\newtheorem{corollary}[theorem]{Corollary}

\theoremstyle{definition}
\newtheorem{definition}[theorem]{Definition}
\newtheorem{example}[theorem]{Example}
\newtheorem{assumption}[theorem]{Assumption}

\theoremstyle{remark}
\newtheorem{remark}[theorem]{Remark}

% --- sets and cones ----------------------------------------------------------
\newcommand{\R}{\mathbb{R}}
\newcommand{\N}{\mathbb{N}}
\newcommand{\Sym}{\mathbb{S}}                      % symmetric matrices
\newcommand{\PSD}{\succeq 0}
\newcommand{\SOS}{\Sigma}                          % the SOS cone, \SOS[x]
\newcommand{\Rx}{\R[x]}                            % polynomial ring
\newcommand{\Mset}{\mathcal{M}}                    % quadratic module

% --- the objects of the certificate ------------------------------------------
\newcommand{\Vdot}{\dot{V}}
\newcommand{\Lie}[2]{\nabla #1 \cdot #2}           % \Lie{V}{f}
\newcommand{\sublevel}[2]{\Omega_{#2}}             % \sublevel{V}{\rho}
\newcommand{\ROA}{\mathcal{R}}                     % region of attraction
\newcommand{\gram}{Q}
\newcommand{\basis}{z}

% --- operators ---------------------------------------------------------------
\DeclareMathOperator{\tr}{tr}
\DeclareMathOperator{\rank}{rank}
\DeclareMathOperator{\diag}{diag}
\DeclareMathOperator{\deg}{deg}
\DeclareMathOperator*{\argmax}{arg\,max}
\DeclareMathOperator*{\argmin}{arg\,min}

% --- a macro for pointing prose at the implementation -------------------------
% Use \codelink{lie_derivative}{julia/src/polynomial_utils.jl} in the text every
% time a symbol first appears. A derivation that names the function implementing
% it is checkable; one that does not is a separate document that happens to sit
% in the same repository.
\newcommand{\codelink}[2]{\texttt{#1} (\texttt{#2})}


\title{Sum-of-Squares Decomposition and the Gram Matrix Parameterisation}
\author{Ali-Eimaan}
\date{}

\begin{document}
\maketitle

\begin{abstract}
% TODO(impl): one paragraph. Checking global nonnegativity of a polynomial is
% NP-hard; checking membership in the SOS cone is a semidefinite program. State
% the trade --- tractability bought at the price of a sufficient-only condition ---
% and say that this document derives the parameterisation and works one example
% end to end, by hand, against the output of \codelink{sos_decomposition}{julia/src/sos_lyapunov.jl}.
\end{abstract}

\section{Setting}
% TODO(impl):
%  - polynomial ring R[x], total degree, the cone of nonnegative polynomials P_{n,2d};
%  - the SOS cone Sigma[x] and its definition;
%  - Sigma subseteq P, with equality only in Hilbert's three cases (n=1;
%    2d=2; n=2, 2d=4). State the classification precisely --- getting it slightly
%    wrong is a tell, and it is easy to state correctly.

\section{The Gram matrix parameterisation}
\begin{proposition}[Gram representation]
% TODO(impl): p of degree 2d is SOS iff there exists Q psd with p = z(x)' Q z(x),
% where z(x) is the vector of monomials of degree at most d.
\end{proposition}
\begin{proof}
% TODO(impl): both directions.
%  (<=) Cholesky Q = L L' gives p = sum_i (L' z)_i^2 explicitly. Write the
%       explicit squares --- this is the direction that produces the certificate
%       and the one the code uses.
%  (=>) Every square in a representation of p has degree at most d, so each is a
%       linear combination of z; collect the coefficient vectors into L.
%       Say why "degree at most d" holds --- it is the Newton polytope argument
%       and it is worth one honest sentence rather than being asserted.
\end{proof}

\begin{remark}[Why this is an SDP]
% TODO(impl): the map Q -> coefficients of z'Qz is linear, so matching
% coefficients with p gives affine equality constraints on Q. Combined with
% Q psd the feasible set is a spectrahedron. Count: Q is N x N with
% N = binom(n+d, d), and there are binom(n+2d, 2d) coefficient equalities, so
% the representation is non-unique whenever N(N+1)/2 exceeds that count. That
% non-uniqueness is exactly the freedom the SDP searches over --- make this point,
% it is the one that turns the parameterisation from a fact into a method.
\end{remark}

\section{Newton polytope reduction}
% TODO(impl): the monomial basis can be restricted to lattice points in half the
% Newton polytope of p (Reznick). State it, give the two-line reason, and report
% the actual size reduction it achieves on the degree-6 Van der Pol Lyapunov
% candidate --- a number from gram_dimension, not an estimate.
% If the implementation does not exploit this, say so here.

\section{A fully worked example}
% TODO(impl): THE SECTION THAT MATTERS. Take a specific quartic in two
% variables, small enough to do entirely by hand, and:
%   1. write z(x) = [1, x1, x2, x1^2, x1 x2, x2^2]';
%   2. write the coefficient-matching equations explicitly, including the
%      free parameter(s) that make Q non-unique;
%   3. exhibit one psd choice of Q;
%   4. take its Cholesky factor and write out the explicit sum of squares;
%   5. verify by expanding, by hand, back to p.
%
% Then print the Q that julia/src/sos_lyapunov.jl returns for the same p and
% compare. It will differ --- the SDP picks an interior point, not your choice ---
% and explaining WHY the two differ while both being valid is a better
% demonstration of understanding than either matrix alone.
%
% Suggested p: something whose hand-decomposition is not the obvious one, so the
% free parameter is genuinely visible. Do not pick a perfect square.

\section{Nonnegative but not SOS}
% TODO(impl): the Motzkin polynomial
%   M(x,y,z) = x^2 y^4 + x^4 y^2 - 3 x^2 y^2 z^2 + z^6.
%  - prove nonnegativity by AM--GM on the three terms;
%  - show it is not SOS. The standard argument: assume M = sum q_i^2, use the
%    Newton polytope to restrict which monomials can appear in the q_i, and
%    derive a contradiction from the coefficient of x^2 y^2 z^2. Write it out;
%    it is short and it is the only place in this document where something is
%    proved to be impossible rather than possible.
%  - report that \codelink{is_sos}{julia/src/sos_lyapunov.jl} returns false for it
%    and name the test that asserts so.

\section{Numerical reality}
% TODO(impl):
%  - the solver returns a floating-point Q; the SOS identity holds to residual,
%    not exactly. Give the residual actually observed, from gram_residual.
%  - conditioning of the monomial basis at high degree, with a measured condition
%    number rather than an adjective;
%  - eigenvalue clipping: what the code does, at what threshold, and what it
%    refuses to do (clip a genuinely negative eigenvalue and call it psd);
%  - rational rounding to obtain an exact certificate --- cite the approach, state
%    clearly that it is NOT implemented here. An honest omission is stronger than
%    an implied capability.

\section{What this buys and what it does not}
% TODO(impl): short and blunt.
%  - SOS feasible => nonnegative: a proof, modulo the numerical caveats above.
%  - SOS infeasible => nothing. Not "no certificate exists": no certificate exists
%    AT THIS DEGREE, in THIS basis, for THIS multiplier structure.
% This section is three sentences and it is why a reader will trust the rest.

\begin{thebibliography}{9}
% TODO(impl): Parrilo (2000); Reznick (1978) on Newton polytopes; Motzkin (1967);
% Blekherman--Parrilo--Thomas (2013); Powers--Wormann on rational certificates.
\end{thebibliography}

\end{document}
